\chapter{Number theory}

\section{Modular arithmetic}
	\kactlimport{ModularArithmetic.h}
	\kactlimport{ModInverse.h}
	\kactlimport{ModPow.h}
	\kactlimport{ModLog.h}
	\kactlimport{ModSum.h}
	\kactlimport{ModMulLL.h}
	\kactlimport{ModSqrt.h}
	\kactlimport{Jacobi.h}

\section{Primality}
	\kactlimport{FastEratosthenes.h}
	\kactlimport{MillerRabin.h}
	\kactlimport{Factor.h}

\section{Divisibility}
	\kactlimport{euclid.h}
	% \kactlimport{Euclid.java}
	\kactlimport{CRT.h}
	\kactlimport{Frobenius.h}

	\subsection{Bézout's identity}
	For $a \neq $, $b \neq 0$, then $d=gcd(a,b)$ is the smallest positive integer for which there are integer solutions to
	$$ax+by=d$$
	If $(x,y)$ is one solution, then all solutions are given by
	$$\left(x+\frac{kb}{\gcd(a,b)}, y-\frac{ka}{\gcd(a,b)}\right), \quad k\in\mathbb{Z}$$

	\kactlimport{phiFunction.h}

\section{Hensel lifting}
	Let $f \in \mathbb Z[x]$, $p$ prime, $k \ge 1$ and $f(r) \equiv 0 \pmod{p^k}$.
	If $f'(r) \not\equiv 0 \pmod p$, then $r' = r - f(r)u$ with $u \equiv f'(r)^{-1} \pmod{p^k}$
	is the unique root mod $p^{2k}$ with $r' \equiv r \pmod{p^k}$. Start from the roots
	mod $p$ and double $k$ each step. If $f'(r) \equiv 0 \pmod p$, then either all of
	$r + tp^k$, $0 \le t < p$, are roots mod $p^{k+1}$ (if $f(r) \equiv 0 \pmod{p^{k+1}}$), or none is.

\section{Vieta jumping}
	For an equation symmetric in $a,b$ and quadratic in each, e.g.\ $a^2+b^2 = k(ab+1)$,
	take a solution with $a \ge b \ge 0$. As a quadratic in $a$ the roots sum to $kb$ and
	multiply to $b^2-k$, so $(a', b)$ with $a' = kb-a = (b^2-k)/a$ is also a solution.
	Keep replacing the larger variable by this value while it stays positive and smaller;
	the descent ends at a minimal solution that must be a base case (here $b=0$, so $k=a^2$).
	Hence $k$ is a perfect square whenever $(a^2+b^2)/(ab+1) = k \in \mathbb N$.

\section{Fractions}
	\kactlimport{ContinuedFractions.h}
	\kactlimport{FracBinarySearch.h}

\section{Pythagorean Triples}
 The Pythagorean triples are uniquely generated by
 \[ a=k\cdot (m^{2}-n^{2}),\ \,b=k\cdot (2mn),\ \,c=k\cdot (m^{2}+n^{2}), \]
 with $m > n > 0$, $k > 0$, $m \bot n$, and either $m$ or $n$ even.

\section{Primes}
	$p=962592769$ is such that $2^{21} \mid p-1$, which may be useful. For hashing
	use 970592641 (31-bit number), 31443539979727 (45-bit), 3006703054056749
	(52-bit). There are 78498 primes less than 1\,000\,000.

	Primitive roots exist modulo any prime power $p^a$, except for $p = 2, a > 2$, and there are $\phi(\phi(p^a))$ many.
	For $p = 2, a > 2$, the group $\mathbb Z_{2^a}^\times$ is instead isomorphic to $\mathbb Z_2 \times \mathbb Z_{2^{a-2}}$.

\section{Estimates}
	$\sum_{d|n} d = O(n \log \log n)$.

	The number of divisors of $n$ is at most around 100 for $n < 5e4$, 500 for $n < 1e7$, 2000 for $n < 1e10$, 200\,000 for $n < 1e19$.

\section{Mobius Function}
\[
	\mu(n) = \begin{cases} 0 & n \textrm{ is not square free}\\ 1 & n \textrm{ has even number of prime factors}\\ -1 & n \textrm{ has odd number of prime factors}\\\end{cases}
\]
  Mobius Inversion:
  \[ g(n) = \sum_{d|n} f(d) \Leftrightarrow f(n) = \sum_{d|n} \mu(d)g(n/d) \]
  Other useful formulas/forms:

  $ \sum_{d | n} \mu(d) = [ n = 1] $ (very useful)

  $ g(n) = \sum_{n|d} f(d) \Leftrightarrow f(n) = \sum_{n|d} \mu(d/n)g(d)$

 $ g(n) = \sum_{1 \leq m \leq n} f(\left\lfloor\frac{n}{m}\right \rfloor ) \Leftrightarrow f(n) = \sum_{1\leq m\leq n} \mu(m)g(\left\lfloor\frac{n}{m}\right\rfloor)$
